\documentclass[12pt,a4paper]{article}
% ============================================================
%  Оформление отчёта. Обычно здесь ничего менять не нужно.
%  Собирается pdfLaTeX (по умолчанию, в т. ч. в Overleaf),
%  XeLaTeX или LuaLaTeX.
% ============================================================

\usepackage{iftex}
\ifPDFTeX
  \usepackage[T2A]{fontenc}
  \usepackage[utf8]{inputenc}
  \IfFileExists{tempora.sty}{\usepackage{tempora}}{}
\else
  \usepackage{fontspec}
  \IfFontExistsTF{Times New Roman}
    {\setmainfont{Times New Roman}}
    {\IfFontExistsTF{CMU Serif}{\setmainfont{CMU Serif}}{\setmainfont{Liberation Serif}}}
  \IfFontExistsTF{Consolas}
    {\setmonofont{Consolas}[Scale=MatchLowercase]}
    {\IfFontExistsTF{Courier New}
      {\setmonofont{Courier New}}
      {\IfFontExistsTF{CMU Typewriter Text}{\setmonofont{CMU Typewriter Text}}{\setmonofont{Liberation Mono}[Scale=MatchLowercase]}}}
\fi
\usepackage[english,russian]{babel}

% --- Поля и абзацы (как в исходном doc-шаблоне курса) ---
\usepackage[a4paper, margin=2.54cm]{geometry}
\setlength{\parindent}{0pt}
\setlength{\parskip}{4pt plus 1pt}
\frenchspacing

% --- Заголовки разделов без номеров ---
\usepackage{titlesec}
\setcounter{secnumdepth}{0}
\titleformat{\section}{\large\bfseries}{}{0pt}{}
\titleformat{\subsection}{\normalsize\bfseries}{}{0pt}{}
\titlespacing*{\section}{0pt}{14pt plus 4pt minus 2pt}{4pt plus 2pt}
\titlespacing*{\subsection}{0pt}{10pt plus 2pt minus 2pt}{3pt plus 1pt}

% --- Таблицы, рисунки, подписи ---
\usepackage[table]{xcolor}
\usepackage{graphicx}
\usepackage{tabularx}
\usepackage{float}
\usepackage{caption}
\DeclareCaptionLabelSeparator{spacedendash}{~-- }
\captionsetup{labelsep=spacedendash, font=small, singlelinecheck=false, justification=raggedright}
\captionsetup[figure]{justification=centering, singlelinecheck=true}
\captionsetup[table]{position=top, skip=4pt}
\DeclareCaptionType[fileext=lol]{listing}[Листинг][Листинги]
\captionsetup[listing]{position=top, skip=4pt, hypcap=false}
\addto\captionsrussian{%
  \renewcommand{\figurename}{Рисунок}%
  \renewcommand{\tablename}{Таблица}%
}
\newcommand{\thead}[1]{\multicolumn{1}{|c|}{\cellcolor{gray!20}\textbf{#1}}}
% Колонки таблиц без растягивания пробелов: L – резиновая, P{ширина} – фиксированная
\newcolumntype{L}{>{\raggedright\arraybackslash}X}
\newcolumntype{P}[1]{>{\raggedright\arraybackslash}p{#1}}

% --- Код и вывод терминала ---
% code     – фрагмент исходного кода;
% terminal – вывод программы, strace и т. п.;
% \inputterminal{файл} / \inputcode{файл} – то же, но из файла.
\usepackage{fvextra}
\fvset{breaklines=true, breakanywhere=true,
       frame=leftline, framerule=1.5pt, framesep=6pt,
       rulecolor=\color{gray!60}, xleftmargin=4pt}
\DefineVerbatimEnvironment{code}{Verbatim}{fontsize=\small}
\DefineVerbatimEnvironment{terminal}{Verbatim}{fontsize=\footnotesize}
\newcommand{\inputcode}[2][]{\VerbatimInput[fontsize=\small,#1]{#2}}
\newcommand{\inputterminal}[2][]{\VerbatimInput[fontsize=\footnotesize,#1]{#2}}
% Листинг с подписью сверху:
%   \begin{listingblock}{Подпись} \begin{code}...\end{code} \end{listingblock}
\newenvironment{listingblock}[1]
  {\par\medskip\captionof{listing}{#1}\nopagebreak\vspace{-\topsep}\vspace{-\parskip}}
  {\par}

% Приложение: с новой страницы, заголовок «Приложение А. ...»,
% листинги внутри нумеруются А.1, А.2, ...
\newcounter{appendixno}
\newcommand{\appendixsection}[2]{%
  \clearpage
  \section{Приложение #1. #2}%
  \stepcounter{appendixno}%
  \setcounter{listing}{0}%
  \renewcommand{\thelisting}{#1.\arabic{listing}}%
  \renewcommand{\theHlisting}{app\arabic{appendixno}.\arabic{listing}}}

% --- Ссылки ---
\usepackage[hidelinks]{hyperref}
\usepackage{xurl}

% --- Подсказки шаблона: серый курсив, удалить перед сдачей ---
\newcommand{\hint}[1]{\leavevmode{\color{gray}\itshape #1}}
% Рамка-заглушка вместо рисунка
\newcommand{\placeholder}[2][4.5cm]{%
  \fbox{\parbox[c][#1][c]{0.95\linewidth}{\centering\hint{#2}}}}

\newcommand{\ReportStudent}{В. Е. Авраменко}
\newcommand{\ReportGroup}{М8О-204БВ-25}
\newcommand{\ReportTeacher}{Н. А. Гапонов}
\newcommand{\ReportDate}{02.11.2026}
\newcommand{\ReportYear}{2026}
\newcommand{\ReportCourse}{Операционные системы и архитектура компьютера}
\begin{document}
\sloppy
% Титульный лист
\begin{titlepage}
  \centering\large

  МОСКОВСКИЙ АВИАЦИОННЫЙ ИНСТИТУТ\\
  (НАЦИОНАЛЬНЫЙ ИССЛЕДОВАТЕЛЬСКИЙ УНИВЕРСИТЕТ)\par
  \vspace{18pt}
  Институт №8 «Компьютерные науки»\\
  Кафедра 806 «Вычислительная математика и программирование»\par

  \vspace{\stretch{1.2}}

  {\bfseries Отчёт по лабораторной работе №4\\
  Вариант 1\\
  по курсу «\ReportCourse»\par}

  \vspace{\stretch{1.6}}

  \begin{flushright}
    \begin{tabular}{r@{\hspace{0.6em}}l}
      Выполнил:      & \ReportStudent \\[2pt]
      Группа:        & \ReportGroup   \\[2pt]
      Преподаватель: & \ReportTeacher \\[2pt]
      Дата:          & \ReportDate    \\[2pt]
      Оценка:        & \rule[-2pt]{4.2cm}{0.4pt} \\
    \end{tabular}
  \end{flushright}

  \vspace{\stretch{1}}

  Москва, \ReportYear
\end{titlepage}
\setcounter{page}{2}
\section{Цель работы}
Приобрести практические навыки работы с файловой системой и организации обмена данными между процессами через отображаемые в память файлы (File Mapping).

\section{Задание: вариант 1}
В лабораторной работе №4 используется вариант из лабораторной работы №2. Родительский процесс получает первой строкой имя выходного файла и запускает отдельную дочернюю программу. Далее пользователь вводит строки целых чисел типа \texttt{int}. Дочерняя программа вычисляет сумму каждой строки и записывает её в указанный файл. Количество чисел в строке произвольное, в пределах ресурсов ОС. Необходимо обрабатывать системные ошибки.

В отличие от ЛР №2, данные между процессами передаются через файловое отображение \texttt{mmap}, а готовность данных и подтверждение обработки обозначаются сигналами. Программы написаны на C11 для Linux. Основой послужили исходники из каталога \texttt{1/src}; функции ввода и вычисления суммы сохранены с небольшими изменениями проверки входа.

\section{Организация программы}
\begin{table}[H]
\small
\begin{tabularx}{\linewidth}{|P{3.4cm}|L|}
\hline
\thead{Файл} & \thead{Назначение}\\ \hline
\texttt{parent.c} & Создание файла обмена и процесса; публикация строк; ожидание подтверждений и завершения.\\ \hline
\texttt{child.c} & Открытие результата; получение команд; вычисление и запись сумм.\\ \hline
\texttt{mapping.c/.h} & Отображение файла, блокировка и синхронное ожидание сигналов.\\ \hline
\texttt{parent\_input.c/.h} & Буферизованное чтение имени файла и строк, ожидание ввода или сигнала.\\ \hline
\texttt{sum.c/.h} & Разбор чисел, проверка диапазона \texttt{int} и переполнения суммы.\\ \hline
\texttt{status.h} & Коды результатов ввода и разбора.\\ \hline
\end{tabularx}
\end{table}

\subsection{Схема взаимодействия}
\begin{terminal}
Пользователь -> parent -> mmap(MAP_SHARED) -> child -> файл сумм
                  |                           |
                  +-------- SIGUSR1 ---------->|
                  |<------- SIGUSR2 -----------+
                  +-------- waitpid ----------> завершение
\end{terminal}
Первый \texttt{SIGUSR2} означает готовность дочерней программы; последующие --- завершение обработки одной строки. В каждый момент времени передаётся не более одной неподтверждённой команды.

\section{Алгоритм обмена}
\begin{enumerate}
\item Родитель читает имя выходного файла и блокирует \texttt{SIGUSR1}, \texttt{SIGUSR2}, \texttt{SIGCHLD}. Маска наследуется при \texttt{fork} и сохраняется при \texttt{exec}.
\item Через \texttt{mkstemp} создаётся уникальный файл обмена с ограниченными правами доступа. Его имя сразу удаляется вызовом \texttt{unlink}. Открытый файл продолжает существовать до закрытия дескрипторов.
\item После \texttt{fork} дочерний процесс вызывает \texttt{execv} для запуска \texttt{./child}. Имя результата, дескриптор файла обмена и исходный PID родителя передаются аргументами. Само отображение через \texttt{exec} не передаётся; дочерняя программа создаёт его заново.
\item Дочерняя программа открывает выходной файл и отправляет родителю \texttt{SIGUSR2}. Родитель ждёт готовности либо \texttt{SIGCHLD}, означающего изменение состояния потомка, в том числе ошибку запуска. После подтверждения он выводит PID обоих процессов.
\item Родитель читает строку. Вызов \texttt{ftruncate} задаёт размер файла, равный длине строки с завершающим нулём. Через \texttt{mmap} с \texttt{MAP\_SHARED} и \texttt{PROT\_READ | PROT\_WRITE} создаётся отображение; \texttt{memcpy} копирует строку, \texttt{munmap} снимает отображение. Затем \texttt{kill} посылает \texttt{SIGUSR1}.
\item Получив сигнал, ребёнок узнаёт размер через \texttt{fstat} и создаёт отображение только для чтения. Проверяются завершающий ноль и отсутствие нулей внутри команды. Числа разбираются, сумма записывается через \texttt{fprintf}; результат \texttt{fflush} проверяется.
\item Только после \texttt{munmap} ребёнок посылает \texttt{SIGUSR2}. Родитель получает подтверждение и может менять размер файла для следующей строки. Это исключает уменьшение файла во время чтения отображения и связанный с этим риск \texttt{SIGBUS}.
\item Пустая строка или EOF завершает ввод: после последнего подтверждения родитель устанавливает размер файла в ноль и посылает \texttt{SIGUSR1}. Ребёнок распознаёт конец данных без вызова \texttt{mmap} с нулевой длиной и завершает работу. Родитель вызывает \texttt{waitpid} и проверяет статус. Ожидание с \texttt{WUNTRACED} позволяет обнаружить остановленного ребёнка и завершить его вместо бесконечного ожидания.
\end{enumerate}

Обработчики сигналов не выполняют вычисления или ввод-вывод: используются \texttt{sigwaitinfo} и \texttt{sigtimedwait}. Заблокированный сигнал, пришедший до начала ожидания, остаётся ожидающим. Обычные сигналы могут объединяться, поэтому протокол строго чередует команду и подтверждение. \texttt{MAP\_SHARED} обеспечивает видимость данных обоим процессам; отдельный \texttt{msync} для долговременного сохранения временного файла здесь не требуется.

\section{Проверка данных и обработка ошибок}
Каждый элемент строки преобразуется функцией \texttt{strtol} с проверкой \texttt{errno}, конца токена и диапазона \texttt{INT\_MIN..INT\_MAX}. Сумма хранится в \texttt{long}; перед сложением проверяются границы \texttt{LONG\_MIN..LONG\_MAX}. В проверенной среде Linux x86\_64 это позволяет суммировать несколько чисел \texttt{int}, даже если результат выходит за диапазон \texttt{int}.

Строка только из пробельных символов пропускается. Строка с неверным числом не записывается; следующие команды обрабатываются, но итоговый код завершения ненулевой. Пустая строка заканчивает ввод, как в исходной ЛР №2. Встроенные нулевые байты в имени или команде считаются ошибкой.

Проверяются ошибки создания и изменения размера файла, отображения и снятия отображения, запуска процесса, сигналов, записи, закрытия и ожидания. Прерванные ожидания повторяются при \texttt{EINTR}. Ошибка открытия результата или \texttt{execv} приводит к \texttt{SIGCHLD}, поэтому ожидание готовности не зависает. При локальной ошибке родитель принудительно завершает ребёнка и собирает его через \texttt{waitpid}. Дочерний процесс раз в секунду ожидания проверяет, сохранился ли исходный PID родителя.

Выходной файл открывается на перезапись. При ожидании ввода \texttt{pselect} атомарно снимает блокировку \texttt{SIGCHLD} на время ожидания. Обработчик отмечает получение сигнала, в том числе когда ввод стал доступен одновременно с ним. Вызов \texttt{waitid} проверяет, завершился или остановился потомок, сохраняя статус для \texttt{waitpid}. Ввод читается порциями через \texttt{read}, поэтому даже незаконченная строка не мешает обнаружить отказ. В тестах родитель завершился и забрал статус потомка при открытом входном потоке.

\section{Сборка и пример запуска}
Команды сборки и запуска:
\begin{terminal}
make
cd src
./parent
\end{terminal}
Введите имя выходного файла, затем строки с числами. Для завершения введите пустую строку или нажмите \texttt{Ctrl+D}. Запуск выполняется из \texttt{src}, поскольку дочерняя программа вызывается по пути \texttt{./child}.

В логе проверочного запуска видны PID обоих процессов и содержимое файла с суммами:
\inputterminal{inc/demo.txt}

\section{Результаты проверки}
Проверки выполнены в Ubuntu / WSL2, ядро Linux 6.6.87.2-microsoft-standard-WSL2, GCC 13.3.0. Сборка с \texttt{-std=c11 -Wall -Wextra -Wpedantic -Werror} прошла без предупреждений. Интеграционный набор содержит 24 успешных сценария; подробности сохранены в \texttt{inc/test-results.json}.

\begin{table}[H]
\small
\begin{tabularx}{\linewidth}{|P{5cm}|L|}
\hline
\thead{Проверка} & \thead{Результат}\\ \hline
Положительные и отрицательные числа & Для \texttt{1 2 3} получено 6, для \texttt{10 -5 7} --- 12.\\ \hline
Границы \texttt{int} & Два \texttt{2147483647} дают \texttt{4294967294}; отрицательная сумма также корректна.\\ \hline
Окончание ввода и пробелы & Проверены EOF, пустая строка, CRLF, отсутствие последнего перевода строки, пробельные команды.\\ \hline
Неверные данные & Буквы, суффиксы, выход за диапазон \texttt{int} и NUL обнаруживаются; неверная сумма не записывается.\\ \hline
Длина и последовательность & 100\,000 чисел в одной строке; 1\,500 команд с изменением размера файла --- успешно.\\ \hline
Ошибки файлов & Отсутствующий каталог и ошибка записи в \texttt{/dev/full} дают ненулевой код без зависания.\\ \hline
Отказы процессов & Отсутствие \texttt{child}, завершение родителя и отказ потомка, в том числе посреди ввода, проверены.\\ \hline
Ресурсы и IPC & Временные имена не остаются; трассировка содержит \texttt{MAP\_SHARED}, сигналы и не содержит \texttt{pipe/pipe2}.\\ \hline
\end{tabularx}
\end{table}

Дополнительно выполнены 23 сценария в сборке с AddressSanitizer и UndefinedBehaviorSanitizer; ошибок не обнаружено. Проверка под \texttt{strace} выполнялась отдельно от санитарной сборки. Результаты: \texttt{inc/sanitizer-results.json}.

Полная трассировка находится в \texttt{inc/ipc-trace.txt}. Ниже показаны характерные операции; номера дескрипторов и PID зависят от запуска:
\begin{terminal}
ftruncate(mapping_fd, length + 1)
mmap(..., PROT_READ|PROT_WRITE, MAP_SHARED, mapping_fd, 0)
munmap(...)
kill(child_pid, SIGUSR1)
mmap(..., PROT_READ, MAP_SHARED, mapping_fd, 0)
munmap(...)
kill(parent_pid, SIGUSR2)
\end{terminal}
Это сокращённая схема по трассе, а не дословный вывод \texttt{strace}. Временный файл, оставшийся открытым после удаления имени, виден в трассе с пометкой \texttt{(deleted)}.

\section{Вывод}
Реализованы две отдельные программы на C, передающие строки целых чисел через файловое отображение. Сигналы задают порядок чтения и записи; подтверждение после обработки предотвращает перезапись ещё используемых данных. Дочерний процесс вычисляет суммы и сохраняет их в файл. Проверены нормальная работа, длинные команды и основные ошибки ввода-вывода и жизненного цикла процессов.

\section{Источники}
\begin{enumerate}
\item Методические указания «Лабораторные работы №4» и «Лабораторные работы №2», предоставленные для задания.
\item Linux man-pages. mmap(2). \url{https://man7.org/linux/man-pages/man2/mmap.2.html}.
\item Linux man-pages. signal(7). \url{https://man7.org/linux/man-pages/man7/signal.7.html}.
\item Linux man-pages. sigwaitinfo(2). \url{https://man7.org/linux/man-pages/man2/sigwaitinfo.2.html}.
\end{enumerate}

\appendixsection{А}{Ключевые фрагменты кода}
\subsection{Передача строки через файловое отображение}
\inputcode[firstline=79,lastline=90]{../src/mapping.c}
\subsection{Ожидание подтверждения или отказа}
\inputcode[firstline=28,lastline=49]{../src/mapping.c}
\subsection{Ожидание ввода с обработкой \texttt{SIGCHLD}}
\inputcode[firstline=57,lastline=77]{../src/parent_input.c}
\end{document}
